\scanpage{28}
% letterhead
U N I V E R S I T E\\
DES SCIENCES ET TECHNIQUES\\
\uline{DU LANGUEDOC}

Montpellier, le 13 mai 1981\\
\uline{Service de la Scolarité}\\
Diffusion : Liste A\\
n\textsuperscript{o} 2724\\
AM/MD

\uline{AVIS DE SOUTENANCE D'UNE THESE DE DOCTEUR-INGENIEUR}

PHYSIOLOGIE DE LA NUTRITION DES VEGETAUX

% printed form
\uline{Nom et prénom du candidat} : THIBAUD Jean Baptiste

\uline{Titre de la thèse} : Mécanisme de l'absorption du nitrate par la racine de maïs.

\uline{Date et lieu de soutenance prévus} : 22 mai 1981 à 14 heures, salle de cours n\textsuperscript{o} 7 - U.S.T.L.

\uline{Nom du directeur de thèse} : Monsieur Claude GRIGNON

\uline{Désignation du laboratoire où la thèse est préparée} : Laboratoire de biochimie et physiologie végétale - Centre de recherches agronomiques du Midi 34060 MONTPELLIER CEDEX.

\uline{Résumé de l'ouvrage} : Le mémoire (81 p.) est introduit par une analyse historique dégageant les données acquises sur le problème de l'absorption du nitrate par les végétaux. Le chapitre suivant présente les méthodes, en particulier la réalisation d'un poste d'électrophysiologie. Quatre chapitres de résultats décrivent respectivement le caractère inductible du transporteur de nitrate, les relations avec les échanges de protons, les propriétés stoechiométriques du système, et ses relations avec la pompe à protons de la membrane plasmique. Le chapitre final de discussion présente une analyse critique du modèle retenu (échange 2 nitrate - 1 hydroxyle), examine ses contraintes énergétiques, et propose des directions de recherche. La liste bibliographique comprend 117 références.

\uline{Durée approximative des recherches ayant abouti à la thèse} : 3 ans

\uline{Composition du jury} : M. C. GRIGNON\\
M. J. GUERN\\
M. \unsure{W.|O.} ULLRICH\\
Mme N. PARIS\\
M. JACOB J.L.

P. Le Secrétaire Général,\\
L'Attachée d'Administration Universitaire,\\
\handwritten{\uline{\unsure{Gentil|A. Gentil}}}\\
\uline{A. GENTIL}
% archive sheet number: 27
